\documentclass[10pt,a4paper]{amsart}
\usepackage[margin=19mm]{geometry}
\usepackage{amsmath,amssymb,mathtools}
\usepackage[scr=rsfs,cal=euler]{mathalfa}
\usepackage{xfrac}
\usepackage[unicode,hidelinks,bookmarks=false]{hyperref}
\newcommand{\ie}{i.e.\ }
\newcommand{\cf}{cf.\ }
\newcommand{\resp}{respectively\ }
\newcommand{\Subsection}{\subsection}
\providecommand{\nobreakdash}{\nobreak\hbox{-}}
\setlength{\emergencystretch}{2em}
\multlinegap0pt
\usepackage{bm}
\usepackage{bbm}
\usepackage{dsfont}
\usepackage{xspace}
\usepackage{cancel}
\usepackage{mathtools}
\usepackage{amsthm}
\makeatletter
\@ifundefined{dfn}{\newtheorem{dfn}{Dfn}}{}
\@ifundefined{prop}{\newtheorem{prop}[dfn]{Proposition}}{}
\makeatother
\newcommand{\Cc}{\mathord{\mathcal{C}^{\infty}_{c}}}
\newcommand{\C}{\mathord{\mathcal{C}^{\infty}}}
\newcommand{\RR}{\mathbb{R}}
\newcommand{\cE}{\mathcal{E}}
\newcommand{\com}{\mathbin{{\scriptstyle \circ }}}
\title[Reflexivity of the space of transversal distributions]{Reflexivity of the space of transversal distributions}
\author{Jure Kali\v{s}nik}
\date{Source: arXiv:2211.07257v1 (CC BY 4.0)}
\begin{document}
\maketitle

\noindent\emph{Source and licence.} Reflexivity of the space of transversal distributions. Version arXiv:2211.07257v1 (CC BY 4.0), distributed under the Creative Commons Attribution 4.0 International licence, as recorded in the arXiv licence metadata for this exact version and shown on the abstract page https://arxiv.org/abs/2211.07257v1. Full rights notice: licenses/arxiv-2211.07257-CC-BY-4.0.txt.\par\smallskip

\noindent\textit{Editorial scope.} Counted unit: the Prop whose source label is \texttt{None} in the article (source0.tex lines 574--580, source label \texttt{None}), delivered together with the article's own complete original proof of it (source0.tex lines 581--632, one \texttt{proof} environment of 52 lines). No mathematics is written, repaired or re-derived: statement and proof are the article's own text line for line. Required context retained uncounted: none. Every definition, display and label of those lines is retained unchanged. Omitted from this package and not redistributed: 1--573, 633--938. Nothing inside a retained excerpt is omitted or altered: every retained body is delivered exactly as the article's lines are written, so no omission receipt of any kind accompanies this package. Citations: the retained excerpts carry no citation command at all, so no bibliography entry of the article is transcribed and no reference is redistributed. Reference adaptation: none was needed and none was made; every reference inside the delivered unit resolves inside it, so no unresolved reference or citation is produced. Printed slips kept literal: no printed word, symbol, hypothesis or number of the counted statement or of its proof was altered. Layout and class: the article's own class is \texttt{amsart} with packages including amsmath, amssymb, amsthm, amscd, xy, mathrsfs, stmaryrd, mathabx; the wrapper is the lane's pinned \texttt{amsart} editorial wrapper at 10pt on A4, which restates the theorem-like environments the retained text needs and adds the article's own macro definitions that the retained text needs (\texttt{\textbackslash C}, \texttt{\textbackslash Cc}, \texttt{\textbackslash RR}, \texttt{\textbackslash cE}, \texttt{\textbackslash com}), with extra wrapper packages (bm, bbm, dsfont, xspace, cancel, mathtools). The delivered numbering is the unit's own. Assets: no retained span contains a figure environment, a graphics-inclusion command, a PDF-page inclusion, a table or a TikZ picture; no image file is copied into this package, the package declares no asset dependency and no image file is redistributed with it. Licence: version arXiv:2211.07257v1 of this article is distributed under the Creative Commons Attribution 4.0 International licence, as recorded in the arXiv licence metadata for this exact version and shown on the abstract page https://arxiv.org/abs/2211.07257v1. A full rights notice is delivered at licenses/arxiv-2211.07257-CC-BY-4.0.txt. Characters: every retained excerpt is pure ASCII, so the delivered engine, XeLaTeX with its default Latin Modern text font, reports no missing character.
\begin{prop}
Let $\pi:\RR^{l}\times N\to \RR^{l}$ be the projection onto $\RR^{l}$. The map 
\[
\Phi:\cE'_{\pi}(\RR^{l}\times N)\to\Cc(\RR^{l},\cE'(N))
\]
is then an isomorphism of locally convex spaces.
\end{prop}

\begin{proof}
We will use the isomorphism $\C(\RR^{l}\times N)\cong\C(\RR^{l},\C(N))$ under which
a function $F\in\C(\RR^{l}\times N)$ is identified with the smooth $\C(N)$-valued function on $\RR^{l}$, given by $x\mapsto F_{x}$,
where $F_{x}=F|\{x\}\times N$. This will enable us to compute partial derivatives of such functions in the directions on the base manifold.

First we show that $\Phi$ is continuous. Take any basic neighbourhood $V_{\textbf{B},\textbf{m},\textbf{e}}$ of zero
in $\Cc(\RR^{l},\cE'(N))$. We can then find a sequence $(\mu_{n})$ of positive real numbers, such that $\bigcup_{n=1}^{\infty}\mu_{n}B_{n}$
is a bounded subset of $\C(N)$. If we denote by $\pi_{N}:\RR^{l}\times N\to N$ the projection onto $N$, the set 
$B=\pi_{N}^{*}\left(\bigcup_{n=1}^{\infty}\mu_{n}B_{n}\right)$ is then a bounded subset of $\C(\RR^{l}\times N)$. Any function
in $B$ is of the form $F\com\pi_{N}$ for some $F\in\bigcup_{n=1}^{\infty}\mu_{n}B_{n}$ and hence corresponds to a
constant map $x\mapsto F$ in $\C(\RR^{l},\C(N))$. For any $u\in\Cc(\RR^{l},\cE'(N))$ and any multi-index $\alpha$ we thus have
\[
D^{\alpha}(u(F\com\pi_{N}))(x)=D^{\alpha}(x\mapsto u_{x}(F))(x)=(D^{\alpha}u)_{x}(F).
\]
Now define a sequence $\textbf{e}'$ by $\epsilon_{n}'=\frac{\mu_{n}\epsilon_{n}}{2}$ and take any $T\in K(B,V_{\textbf{m},\textbf{e}'})$.
For any $x\in K_{n-1}^{c}$ and any multi-index $\alpha$ with $|\alpha|\leq m_{n}$ we have
\begin{align*}
p_{B_{n}}((D^{\alpha}\Phi(T))_{x})&=\sup_{F\in B_{n}}|(D^{\alpha}\Phi(T))_{x}(F)|
=\sup_{F\in B_{n}}|D^{\alpha}(T(F\com\pi_{N}))(x)|, \\
&=\frac{1}{\mu_{n}}\sup_{F\in \mu_{n}B_{n}}|D^{\alpha}(T(F\com\pi_{N}))(x)|, \\
&\leq \frac{1}{\mu_{n}}\sup_{G\in B}|D^{\alpha}(T(G))(x)|.
\end{align*}
Since $G\in B$ and $T\in K(B,V_{\textbf{m},\textbf{e}'})$, it follows that $T(G)\in V_{\textbf{m},\textbf{e}'}$. So for 
$x\in K_{n-1}^{c}$ and any multi-index $\alpha$ with $|\alpha|\leq m_{n}$ we have $|D^{\alpha}(T(G))(x)|<\epsilon_{n}'$, which shows that
\[
p_{B_{n}}((D^{\alpha}\Phi(T))_{x})\leq \frac{1}{\mu_{n}}\sup_{G\in B}|D^{\alpha}(T(G))(x)|<\epsilon_{n}
\] 
and hence $\Phi(K(B,V_{\textbf{m},\textbf{e}'}))\in V_{\textbf{B},\textbf{m},\textbf{e}}$.

To see that $\Phi$ is an open map, take any basic neighbourhood $K(B,V_{\textbf{m},\textbf{e}})$ of zero
in $\cE'_{\pi}(\RR^{l}\times N)$. Now define a sequence $(B_{n})$ of bounded subsets of $\C(N)$ by
\[
B_{n}=\{(D^{\gamma}F)_{x}\,|\,x\in K_{n},|\gamma|\leq m_{n},F\in B\}.
\]
Furthermore, define a sequence $\textbf{e}'$ by $\epsilon_{n}'=\frac{\epsilon_{n}}{2^{m_{n}}}$. For any 
$u\in V_{\textbf{B},\textbf{m},\textbf{e}'}$, any $x\in K_{n-1}^{c}$, any $F\in B$ and any multi-index $\alpha$ 
with $|\alpha|\leq m_{n}$ we then have
\begin{align*}
|D^{\alpha}(u(F))(x)|&=|D^{\alpha}(x\mapsto u_{x}(F_{x}))(x)|
=\left|\sum_{\alpha=\beta+\gamma}(D^{\beta}u)_{x}(D^{\gamma}F)_{x}\right|, \\
&\leq \sum_{\alpha=\beta+\gamma}|(D^{\beta}u)_{x}(D^{\gamma}F)_{x}|.
\end{align*}
By definition we have $(D^{\gamma}F)_{x}\in B_{n'}$ for some $n'\geq n$, so we have a bound 
\[
|(D^{\beta}u)_{x}(D^{\gamma}F)_{x}|\leq p_{B_{n'}}((D^{\beta}u)_{x})<\epsilon_{n'}'<\epsilon_{n}'.
\]
Since there are at most $2^{m_{n}}$ possibilities of writing $\alpha=\beta+\gamma$, we get that
\[
|D^{\alpha}(u(F))(x)|<\epsilon_{n}.
\]
This shows that $V_{\textbf{B},\textbf{m},\textbf{e}'}\subset\Phi(K(B,V_{\textbf{m},\textbf{e}}))$ which concludes the proof.
\end{proof}

\end{document}
